\section{Storage codes on sparse side information}\label{sec:codes}

The second hypothesis of Theorem~D is an all-codes index-coding statement, the crux (\cref{def:crux}): on the side information
produced by a bounded-depth genie, every index code for the transposition demand must broadcast a constant fraction of
the input. An oblivious linear-time transposer whose dependency DAG has cheap genies of bounded depth needs the
opposite, broadcasts of $\eps L$ bits with $\eps$ as small as the trade-off of \cref{prop:tradeoff} allows. Rate~$\tfrac12$ is the natural threshold: when the receivers can be paired so
that each knows its partner's demand, broadcasting the XOR of each pair serves everyone with $L/2$ bits, and no
construction we know does better on bounded-degree side information of growing odd girth.

In the hardest instances we know, coming from schedules whose head motion expands (the ones \cref{sec:motion} leaves
open), the side information that matters is symmetric and sparse: two receivers that know each other's demands form an
edge of a bounded-degree graph that looks locally like a tree. This is a heuristic description, not a theorem, but it
singles out the following question, the undirected form of the crux. By the duality between index codes and storage codes
(\cref{sec:prelim-codes}), it asks whether storage codes can beat rate~$\tfrac12$.

\begin{question}[Undirected crux]\label{prob:undirected-crux}
Let $(G_n)$ be graphs of maximum degree at most $\Delta$ that are \emph{locally bipartite}: for each fixed $\ell$, only
$o(\card{V(G_n)})$ vertices lie on odd cycles of length at most $\ell$. Does every storage code on $G_n$ have rate at most
$\tfrac12+o(1)$, over every alphabet?
\end{question}

A negative answer is necessary for the kind of coding gain a linear-time transposer would exploit on such side
information; a positive answer, in the directed and depth-bounded generality of \cref{def:crux}, would combine with depth
reduction to refute oblivious~(T). This section proves the positive answer for restricted classes of codes and of
graphs, and shows where the remaining difficulty lies:
\begin{itemize}[leftmargin=2em]
\item every storage code, linear or not, on any graph that maps homomorphically to $C_5$, to the circular clique
$K_{14/5}$ or to the Petersen graph has rate at most $\tfrac12$, with no degree assumption (\cref{cor:H2});
\item on triangle-free graphs of maximum degree~$3$, no linear storage code over any field beats rate~$\tfrac12$
(Theorem~S, \cref{thm:S});
\item with bounded complex coefficients, no linear code beats rate $\tfrac12+o(1)$ on locally bipartite graphs of
bounded degree, and over finite fields, with a fixed number of symbols per vertex, a gain must come from $p$-torsion
(Theorem~A, \cref{thm:A,cor:torsion});
\item at degree~$5$ the threshold fails: the Clebsch graph carries a binary linear code of rate $5/8$, and no storage code
does better, so its guessing number is exactly~$10$ (\cref{cor:H4});
\item local arguments inside tree-like regions cannot give rate $\tfrac12+o(1)$ at degree at least~$4$, real eigenvalue
codes cannot reach rate~$\tfrac12$ at girth at least~$7$, and Mycielski's construction does not preserve goodness;
whether goodness bounds the fractional chromatic number is open (\cref{sec:ana-limits,sec:ana-targets}).
\end{itemize}
Every numbered result of \cref{sec:entropic,sec:thmH,sec:symmetrization,sec:targets,sec:clebsch,sec:subcubic} is
formalized, apart from the classical Tutte--Berge formula used in the matching form of Theorem~S (\cref{thm:S-barrier});
\cref{sec:ana-A,sec:ana-fields,sec:ana-limits,sec:ana-targets} are not formalized.

\subsection{Entropic polymatroids}\label{sec:entropic}

\begin{lemma}[Entropy of a recoverable word]\label{lem:entropic}
\formal{StorageCode.lean}{codeEntropy, codeEntropy_isPolymatroid, codeEntropy_isValid}
Let $X\colon\Omega\to\alpha^V$ be a random word on a finite probability space. Then $S\mapsto H(X_S)$ is a polymatroid on
$V$. If $X$ is recoverable for $N$, this polymatroid is valid for~$N$.
\end{lemma}

\begin{outline}
$H(X_S)$ is the entropy of the observation $\omega\mapsto X(\omega)|_S$, written as a word in $(\alpha\cup\{\bot\})^V$ that
is $\bot$ outside $S$, in the second paper's entropy library~\cite{SecondPaper2026}. The empty observation is constant,
so it has entropy~$0$. Monotonicity holds because $X_S$ is a function of $X_T$ for $S\subseteq T$. Submodularity is that
library's strong subadditivity for the observations $X_S$, $X_T$ and $X_{S\cap T}$. Validity: recoverability makes
$X_{N(v)\cup\{v\}}$ a function of $X_{N(v)}$, so the two entropies agree.
\end{outline}

\begin{theorem}[Rate bounds bound codes]\label{thm:entropic-rate}
\formal{StorageCode.lean}{codeEntropy_univ_le_of_rateBound, codeEntropy_univ_uniform, log_card_le_of_rateBound}
Suppose $N$ satisfies the rate bound $c\ge0$. Then every random word $X$ recoverable for $N$ satisfies
$H(X_V)\le c\sum_vH(X_v)$, and every storage code $C$ for $N$ over a nonempty alphabet $\alpha$ satisfies
$\log\card C\le c\,\card V\log\card\alpha$.
\end{theorem}

\begin{outline}
The first claim applies the rate bound to the polymatroid of \cref{lem:entropic}. For the second, take $X$ uniform on
$C$; then $H(X_V)=\log\card C$, and each $H(X_v)\le\log\card\alpha$.
\end{outline}

\paragraph{Why it matters.} This is the bridge from polymatroids to codes: a rate bound proved for valid polymatroids,
by a certificate or by Theorem~H, bounds every storage code, linear or not, over every alphabet. For~(T), a rate
bound~$c$ on the side information of the cones gives the crux at rate~$1-c$ (\cref{rem:dec-crux}(b)).

Valid polymatroids inherit the recovery structure as a closure operation, which every certificate below uses.

\begin{lemma}[Closure]\label{lem:closure}
\formal{Polymatroid.lean}{IsPolymatroid.eq_of_closure, IsPolymatroid.insert_eq_of_subset}
Let $h$ be a polymatroid valid for $N$, and let $S\subseteq T$ be such that $N(v)\subseteq S$ for every
$v\in T\setminus S$. Then $h(T)=h(S)$.
\end{lemma}

\begin{outline}
Submodularity gives diminishing returns: if adding $v$ to $A$ does not increase $h$, then adding $v$ to any $B\supseteq A$
does not either. Apply this with $A=N(v)$ and $B\supseteq S$, adding the vertices of $T\setminus S$ one at a time.
\end{outline}

\subsection{Theorem H: rate bounds pull back along homomorphisms}\label{sec:thmH}

\begin{theorem}[Theorem H]\label{thm:H}
\formal{Polymatroid.lean}{pushRank, IsPolymatroid.pushRank, IsValid.pushRank, pushRank_univ, sum_pushRank_singleton_le, RateBound.of_hom, RateBound.of_graphHom, RateBound.of_digraphHom}
Let $N_G$ and $N_T$ be side information on $V$ and $W$, and let $\varphi\colon V\to W$ map $N_G(u)$ into $N_T(\varphi(u))$
for every $u$. For a polymatroid $h$ on $V$ put $h'(Z)=h(\varphi^{-1}(Z))$. Then $h'$ is a polymatroid on $W$, valid for
$N_T$ if $h$ is valid for $N_G$, with $h'(W)=h(V)$ and $\sum_th'(\{t\})\le\sum_vh(\{v\})$. Consequently, if $N_T$ satisfies
the rate bound $c\ge0$, so does $N_G$. In particular this holds for every graph homomorphism $G\to T$ and every digraph
homomorphism (arcs to arcs).
\end{theorem}

\begin{outline}
Preimages preserve unions, intersections and inclusions, so $h'$ is a polymatroid. Validity: if $u$ lies in the fibre of
$t$, then $N_G(u)\subseteq\varphi^{-1}(N_T(t))$, so by \cref{lem:closure} adding the whole fibre of $t$ to
$\varphi^{-1}(N_T(t))$ does not increase~$h$. No looplessness or injectivity is needed. The fibres partition $V$, and
subadditivity bounds $h$ of each fibre by the sum over its vertices. Finally
$h(V)=h'(W)\le c\sum_th'(\{t\})\le c\sum_vh(\{v\})$.
\end{outline}

The formal theorem is about all polymatroids, so it concerns the Shannon bound. The same argument shows that the
push-forward of an entropic polymatroid is entropic: if $h(S)=H(X_S)$, then $h'(Z)=H(Y_Z)$ for the random word
$Y_t=X_{\varphi^{-1}(t)}$. We do not need this, because every target below is certified by Shannon inequalities.

\paragraph{Why it matters.} Theorem~H turns one finite certificate into a statement about infinitely many graphs: a rate
bound for a fixed target~$T$ holds for every side information that maps to~$T$, whatever its size or degree. This is the
mechanism behind every all-codes result of this section about infinitely many graphs, and
the reason a target's structure, not the size of the side
information, decides goodness. For~(T), it underlies the homomorphism route to the crux (\cref{sec:ana-targets}). The
theorem is folklore-level (the proof is one line); related monotonicity statements
for index coding appear in~\cite{EbrahimiSiavoshani2014}, and blow-up lemmas in~\cite{Baber2017,CameronDangRiis2016}.

\subsection{Symmetrization}\label{sec:symmetrization}

\begin{theorem}[Symmetrization]\label{thm:symmetrization}
\formal{Symmetry.lean}{IsAut, symmetrize, IsPolymatroid.symmetrize, IsValid.symmetrize, symmetrize_invariant, symmetrize_univ, sum_symmetrize_singleton, RateBound.of_invariant}
Let $N$ be side information on $V$. If every valid polymatroid that is invariant under all automorphisms of $N$ satisfies
$h(V)\le c\sum_vh(\{v\})$, then $N$ satisfies the rate bound~$c$.
\end{theorem}

\begin{outline}
Average a valid polymatroid over the finite group of permutations $\sigma$ of $V$ with $\sigma(N(v))=N(\sigma(v))$. Each
relabelled function $S\mapsto h(\sigma S)$ is a valid polymatroid, so the average is one too. It is invariant, because
right multiplication by an automorphism permutes the group, and it has the same value on $V$ and the same sum over
singletons as~$h$.
\end{outline}

\paragraph{Why it matters.} With symmetrization, a certificate (\cref{sec:targets}) needs one elemental inequality per orbit, and each set it
mentions can be replaced by an
orbit representative through an explicit automorphism. This is what makes the Clebsch certificate below small enough to
check.

\subsection{Good targets}\label{sec:targets}

A certificate for a rate bound $c$ is a nonnegative combination of \emph{elemental inequalities}
$h(K\cup\{i\})+h(K\cup\{j\})\ge h(K\cup\{i,j\})+h(K)$ that equals $c\sum_vh(\{v\})-h(V)$ once every set is replaced by its
closure (\cref{lem:closure}). We found the certificates by linear programming, re-solved them exactly in rational
arithmetic on their supports, and checked the Petersen, $K_{14/5}$ and Clebsch certificates in Lean with
\texttt{linear\_combination} and exact coefficients (the five-inequality pentagon certificate is closed by
\texttt{linarith}); the generator is \leanfile{tools/certificates.py}. The smallest example shows the mechanism.

\begin{proposition}[The pentagon is good]\label{prop:C5}
\formal{Pentagon.lean}{pentagonNbhd, nbhd_cycleGraph_five, rateBound_pentagon, rateBound_cycleGraph_five}
Every valid polymatroid on the $5$-cycle $0\text{-}1\text{-}2\text{-}3\text{-}4\text{-}0$ satisfies $2h(V)\le\sum_vh(\{v\})$.
\end{proposition}

\begin{outline}
Writing $I(i;j\mid K)$ for the elemental quantity above, the five non-adjacent inequalities
\[
  I(0;2)+I(0;3\mid1,2)+I(1;3)+I(1;4)+I(2;4\mid1)\;\ge\;0
\]
sum to $\sum_vh(\{v\})-2h(V)$ after the closures $\{0,1,2\}\to\{0,2\}$, $\{1,2,3\}\to\{1,3\}$,
$\{0,1,2,3\}\to V$ and $\{1,2,4\}\to V$. The Lean proof is exactly this computation.
\end{outline}

The \emph{circular clique} $K_{p/q}$ has vertices $\Z/p$, with $i\sim j$ if and only if $q\le(j-i)\bmod p\le p-q$.

\begin{proposition}[Further good targets]\label{prop:targets}
\formal{Petersen.lean}{petersenGraph, petersenGraph_adj_iff, petersenKneser_card, petersenKneser_injective, rateBound_petersenGraph}
\formal{CircularClique.lean}{circularClique, rateBound_circularClique_14_5, circularClique_adj_of_lt, circularCliqueHom, rateBound_circularClique_of_le, cycleToCircularClique, rateBound_cycleGraph_odd}
The Petersen graph (the Kneser graph $K(5,2)$) and the circular clique $K_{14/5}$ (the Andrásfai graph
$\mathrm{And}(5)$) are good. If $p/q\le p'/q'$ with $q'\ge1$, the map $i\mapsto\lfloor p'i/p\rfloor$ is a homomorphism
$K_{p/q}\to K_{p'/q'}$ (Bondy and Hell~\cite{BondyHell1990}). Hence every circular clique $K_{p/q}$ with $p/q\le14/5$ is
good, for example $K_{8/3}$ (the Möbius ladder $\mathrm{And}(3)$) and $K_{11/4}$ ($\mathrm{And}(4)$), and so is every
odd cycle $C_{2k+1}$ with $k\ge2$.
\end{proposition}

\begin{outline}
The Petersen certificate is a
combination of $58$ elemental inequalities with weights $\tfrac14$ and~$\tfrac12$ and $25$ closure equalities (of the $40$ that the
file proves), and the labelling is identified with $K(5,2)$ by an explicit bijection. For $K_{14/5}$ the polymatroid is first symmetrized over the
dihedral group (\cref{thm:symmetrization}); the certificate then uses $32$ orbit rows, $21$ explicit automorphisms, $116$
invariance equalities and $13$ closure equalities. For the homomorphism, if $i<j$ and $q\le j-i\le p-q$,
then $\lfloor p'i/p\rfloor+q'\le\lfloor p'j/p\rfloor\le\lfloor p'i/p\rfloor+p'-q'$ because $pq'\le p'q$, and adjacency is
symmetric. The odd cycle $C_{2k+1}$ maps to $K_{(2k+1)/k}$ by $i\mapsto ki$, which turns the differences $\pm1$ into
$\pm k$, and $(2k+1)/k\le14/5$ for $k\ge2$. Theorem~H applies.
\end{outline}

\begin{corollary}[Corollary H2: rate $\tfrac12$ for every code]\label{cor:H2}
\formal{RateHalf.lean}{card_sq_le_of_rateBound_half, IsGood.card_sq_le, IsGood.codeEntropy_le}
\formal{Polymatroid.lean}{IsGood.of_hom}
Let $G$ be a good graph, for instance a graph with a homomorphism to $C_5$, to the Petersen graph or to a circular clique
$K_{p/q}$ with $p/q\le14/5$. Then every storage code $C\subseteq\alpha^{V}$ on $G$ over a nonempty alphabet, linear or
not, satisfies $\card C^2\le\card\alpha^{\card V}$, and every random word $X$ recoverable for $G$ satisfies
$2H(X_V)\le\sum_vH(X_v)$.
\end{corollary}

\begin{outline}
Goodness pulls back along homomorphisms (Theorem~H); combine with \cref{prop:C5,prop:targets} and
\cref{thm:entropic-rate}. The integer form takes logarithms.
\end{outline}

The graphs with a homomorphism to some $K_{p/q}$ with $p/q\le14/5$ are exactly the graphs of circular chromatic number at
most $14/5$, since for a finite graph the circular chromatic number is the least $p/q$ such that it maps to
$K_{p/q}$~\cite{Vince1988,Zhu2001}.

A graph with a perfect matching carries the pairing code (copy one symbol across each matched edge), of rate exactly
$\tfrac12$, so for good graphs with a perfect matching the corollary determines the guessing number:
$\gn(G)=\card V/2$. The corollary strictly
extends the guessing-number results of Mazumdar, McGregor and Vorotnikova~\cite[Thms~9--11]{MMV2019} (their graphs map to
$K_{8/3}$ or to $C_5$), and contains the odd-cycle theorem of Christofides and Markström~\cite{ChristofidesMarkstrom2011}.
For Petersen it makes rigorous the numerical values of Blasiak, Kleinberg and Lubetzky~\cite{BKL2013}.

\paragraph{Why it matters.} \Cref{cor:H2} is an all-codes answer to the undirected crux, \cref{prob:undirected-crux}, for
every graph that maps to a good target: the graphs need not have bounded degree, and they include graphs of arbitrarily
large odd girth, for example long odd cycles. The question is
whether high-girth bounded-degree graphs map to good targets. For cubic graphs this is a known open problem:

\begin{corollary}[Corollary H3: a conditional degree-$3$ result]\label{cor:H3}
\formal{RateHalf.lean}{IsGood.card_sq_le}
\formal{Polymatroid.lean}{IsGood.of_hom}
\formal{Pentagon.lean}{rateBound_cycleGraph_five}
\formal{CircularClique.lean}{rateBound_circularClique_14_5}
If every cubic graph of sufficiently large girth maps to $C_5$ (Nešetřil's Pentagon Problem~\cite{Nesetril1999}), or to
$K_{14/5}$, then every storage code on such a graph, nonlinear included, has rate at most $\tfrac12$.
\end{corollary}

\begin{outline}
This is \cref{cor:H2} applied under the stated hypothesis.
\end{outline}

Whether high-girth cubic graphs have circular chromatic number below~$3$ is itself open~\cite{DaneshgarMadani2015}. Some
targets are known to fail: random cubic graphs do not map to $C_7$~\cite{Hatami2005}, and the Clebsch graph, to which every
subcubic graph of girth at least~$17$ maps by a theorem of DeVos and Šámal~\cite{DeVosSamal2006}, is not good
(\cref{cor:H4}). For linear codes, Theorem~S below settles degree~$3$ unconditionally.

\subsection{The Clebsch graph: guessing number ten}\label{sec:clebsch}

The Clebsch graph is the folded $5$-cube: vertices $0,\dots,15$, with $i\sim j$ if and only if $i\oplus j$ (bitwise
exclusive or) is a power of two or equals $15$. It is triangle-free and $5$-regular.

\begin{corollary}[Corollary H4: $\gn(\text{Clebsch})=10$]\label{cor:H4}
\formal{ClebschCertificate.lean}{rateBound_clebsch, rateBound_clebschGraph, clebsch_card_le, clebsch_guessingNumber, clebschCode_eq, mem_ker_clebsch_iff, finrank_ker_clebsch, rank_clebsch}
\formal{Clebsch.lean}{clebschGraph, clebschGraph_adj_iff, clebschCode, clebschCode_card, clebschCode_isStorageCode, exists_clebsch_storageCode}
Every valid polymatroid on the Clebsch graph satisfies $h(V)\le\tfrac58\sum_vh(\{v\})$. Hence every storage code on the
Clebsch graph over an alphabet of size $q$ has at most $q^{10}$ codewords, while the binary linear storage code
$\ker(I+A)$ over $\F_2$, where $A$ is the adjacency matrix, has exactly $2^{10}$. So $\gn(\text{Clebsch})=10$, and
$I+A$ has rank~$6$ over $\F_2$.
\end{corollary}

\begin{outline}
Symmetrize over the automorphism group, of order $1920$ (\cref{thm:symmetrization}). The certificate has $54$ orbit rows;
unlike the certificates for $C_5$, the Petersen graph and $K_{14/5}$, it uses elemental inequalities between adjacent
vertices as well. It uses $126$ explicit
automorphisms, $225$ invariance equalities and $54$ closure equalities, and is checked with exact rational coefficients.
For the lower bound, the code is the span of ten explicit vectors, each satisfying $x_v=\sum_{u\sim v}x_u$ over
$\F_2$, in systematic form on ten coordinates; injectivity of the encoding gives $2^{10}$ codewords. Since $\ker(I+A)$ is itself a storage code, the upper bound shows that
it has at most $2^{10}$ words, so it equals the code; its dimension is $10$, and the rank of $I+A$ is $16-10=6$. The
identification with the folded $5$-cube is checked by decision procedure.
\end{outline}

\begin{proposition}[The Shannon bound of the Clebsch graph]\label{prop:clebsch-exact}
\formal{ClebschCertificate.lean}{clebsch_rateBound_iff}
\formal{Polymatroid.lean}{RateBound.mono}
The Clebsch graph satisfies the rate bound $c$ if and only if $c\ge\tfrac58$.
\end{proposition}

\begin{outline}
The \emph{if} direction is \cref{cor:H4} with monotonicity in $c$. For the \emph{only if} direction, the entropic
polymatroid of the uniform distribution on the binary code has $h(V)=10\log2$ and $h(\{v\})\le\log2$, which violates
every rate bound below~$\tfrac58$.
\end{outline}

Cameron, Dang and Riis proved $10\le\gn(\text{Clebsch})\le11$~\cite{CameronDangRiis2016} and, in the preprint
version, asked for the exact guessing numbers of the strongly regular graphs they consider~\cite[Problem~5.11]{CameronDangRiis2014};
\cref{cor:H4} answers this for the Clebsch graph. The binary code $\ker(I+A)$ goes back to Blasiak,
Kleinberg and Lubetzky~\cite{BKL2013}.

\paragraph{Why it matters.} The Clebsch graph is triangle-free with odd girth~$5$ and degree~$5$, and its codes beat the
pairing by a constant factor. So bounded degree alone does not force rate~$\tfrac12$: any proof of the undirected crux must
use local bipartiteness (few short odd cycles), not only sparsity. It also rules out
targets: it is not good, although every subcubic graph of girth at least~$17$ maps to it~\cite{DeVosSamal2006}, and
by Theorem~H no graph that receives it is good, which makes iterated Mycielski graphs bad (\cref{prop:M}). On the
construction side it shows what a coding gain can look like
at small girth: the binary code is the reduction modulo~$2$ of a rational eigenvalue code with coefficients~$1$
(\cref{rem:ana-clebsch}), a mechanism that Theorem~A and \cref{cor:torsion} rule out on locally bipartite families of
bounded degree.

\subsection{Theorem S: linear codes on subcubic triangle-free graphs}\label{sec:subcubic}

Let $F$ be a field and $M\in F^{V\times V}$ a matrix with unit diagonal. The \emph{recovery set} of $v$ is
$S_v=\{u\ne v: M_{vu}\ne0\}$. The matrix has \emph{no recovery triangle} if $q\notin S_p$ whenever $p,q\in S_r$; this
holds whenever the off-diagonal support of $M$ lies in the edges of a triangle-free graph. A vertex satisfies
\emph{reciprocity}~$(\star\star)$ if $\card{S_v}\le2$, or $\card{S_v}=3$ and $r\in S_v$ whenever $v\in S_r$.

\begin{theorem}[Theorem S]\label{thm:S}
\formal{Subcubic.lean}{Star, RecTriangleFree, ThmSHyp, finrank_subCode_le_half, thmSHyp_of_cliqueFree, star_of_degree_le_three, finrank_ker_le_half_of_degree_le_three, rank_ge_half_of_degree_le_three}
Let $M\in F^{V\times V}$ have unit diagonal and no recovery triangle. For every vertex set $W$ on which every vertex
satisfies $(\star\star)$, the kernel of the principal submatrix $M[W,W]$ has dimension at most $\lfloor\card W/2\rfloor$.
In particular, if $G$ is a triangle-free graph of maximum degree at most~$3$ and the off-diagonal support of $M$ lies in
the edges of $G$, then $\dim\ker M\le\lfloor\card V/2\rfloor$ and $\rank M\ge\lceil\card V/2\rceil$, over every field.
\end{theorem}

\begin{outline}
Codes on vertex sets are subspaces of $F^V$: the code on $W$ is the set of vectors supported on $W$ that satisfy the rows of
$W$, so deleting vertices is an inclusion of subspaces. The proof is a strong induction over $W$; suppose $W$ is a minimal
counterexample.
\begin{enumerate}[leftmargin=2em]
\item \emph{Start.} Deleting one vertex costs at most one dimension, so $\card W=2m+1$ is odd, the code has dimension
exactly $m+1$, and every singleton deletion is \emph{tight}: it leaves the largest dimension the theorem allows.
\item \emph{Tightness.} On a tight deletion no coordinate vanishes identically, since otherwise a smaller set would be a
counterexample. If two coordinates $p,q$ vanish together on a tight code, deleting both is again tight, and the deleted
code is the subspace $\{y_p=0\}$ of the old one.
\item \emph{Maximal chains.} A \emph{good chain} starts at a vertex $v_0$ with at most two recovery neighbours in $W$, which
exists by a parity argument (otherwise the recovery relation would be a symmetric $3$-regular graph on an odd set). It
repeatedly adds a pair $\{p,q\}$ that is exactly the part of a chain vertex's recovery set outside the chain. Along a good
chain, every deletion is tight and its code lies inside the code of~$W$, so every row of a chain vertex restricted to
the rest is a relation on the remaining code. If some chain vertex has exactly one recovery neighbour outside the chain,
that coordinate vanishes, contradicting tightness. If one has exactly two, the chain extends. Reciprocity rules out
three, because the vertex that added a chain vertex lies in its recovery set.
\item \emph{Final count.} A maximal chain therefore leaves every chain vertex with no recovery neighbours outside it. Call
\emph{free} the vertex $v_0$ and the first vertex of each pair. A codeword that vanishes on the free vertices vanishes on
the whole chain, because the row of each pair's parent then forces the second vertex of the pair to vanish. For the last
pair $\{p,q\}$ even $p$ is forced: its row involves only earlier chain vertices, since $q\notin S_p$ ($p,q\in S_r$ for the
parent $r$, and $M$ has no recovery triangle). So a codeword is determined, on the chain, by $k$ coordinates for a chain
with $k$ pairs, and $m+1\le k+\lfloor(\card W-2k-1)/2\rfloor=m$, a contradiction.
\end{enumerate}
\end{outline}

The proof needs neither Gallai's lemma nor connectivity, and the first part of \cref{thm:S}, stated for recovery sets, is slightly stronger
than the graph-theoretic statement. The triangle condition is only needed on the recovery sets, and reciprocity only on the vertex
set being bounded. Over a field, linear storage codes are exactly the subspaces of the kernels of such matrices:

\begin{lemma}[Linear storage codes are subspaces of kernels]\label{lem:linear-pattern}
\formal{Subcubic.lean}{IsLinearStorageCode, IsLinearStorageCode.exists_coeffs, IsLinearStorageCode.exists_pattern, isLinearStorageCode_of_le_ker, isStorageCode_iff_isLinearStorageCode}
Call a subspace $C\subseteq F^V$ a \emph{linear storage code} on $G$ if every $x\in C$ that vanishes on the neighbours
of~$v$ vanishes at~$v$; over a finite field this is equivalent, by linearity, to the defining property of
\cref{def:storage-code}. If $C$ is a linear storage code, each coordinate is a fixed linear combination of the
neighbouring coordinates on~$C$, so $C\subseteq\ker M$ for a matrix $M$ with unit diagonal and off-diagonal support in the
edges of~$G$. Conversely, every subspace of such a kernel is a linear storage code.
\end{lemma}

\begin{outline}
On $C$, the evaluation at $v$ vanishes on the common kernel of the evaluations at the neighbours of $v$, so it lies in
their span (Mathlib's \texttt{mem\_span\_of\_iInf\_ker\_le\_ker}). Conversely, on $\ker M$ the row of~$v$ gives
$x_v=-\sum_{u\sim v}M_{vu}x_u$. For the equivalence, $x,y\in C$ agree on a set exactly when $x-y\in C$ vanishes there.
\end{outline}

\begin{corollary}[Theorem S for linear storage codes]\label{cor:S-linear}
\formal{Subcubic.lean}{IsLinearStorageCode.finrank_le_half, IsLinearStorageCode.card_sq_le}
On a triangle-free graph of maximum degree at most~$3$, every linear storage code over any field has dimension at most
$\lfloor\card V/2\rfloor$; over a finite field $F$ it has $\card C^2\le\card F^{\card V}$.
\end{corollary}

\begin{outline}
By \cref{lem:linear-pattern}, $C\subseteq\ker M$ for a matrix $M$ as in \cref{thm:S}; apply \cref{thm:S}. Over a finite
field, $\card C=\card F^{\dim C}$.
\end{outline}

The exact minrank follows by splitting along separated vertex sets.

\begin{theorem}[Theorem S, barrier and matching forms]\label{thm:S-barrier}
\partformal{Subcubic.lean}{finrank_subCode_union_le, finrank_subCode_biUnion_le, finrank_ker_le_barrier, finrank_ker_le_heavy, exists_linearStorageCode_of_matching, exists_pattern_finrank_ker_ge_of_matching}{the identification of the minimum of the bound with the matching number, which is the Tutte--Berge formula~\cite{Berge1958}.}
Let $M$ have unit diagonal, no recovery triangle, and off-diagonal support in the edges of a graph~$G$. Let $U\subseteq V$ be
a set outside which every vertex satisfies $(\star\star)$, and let the vertices outside $U$ be partitioned into blocks with
no edges between different blocks (for instance the components of $G-U$). Then
\[
  \dim\ker M\;\le\;\card U+\sum_{B}\lfloor\card B/2\rfloor .
\]
Taking for $U$ the vertices violating $(\star\star)$ gives $\dim\ker M\le\card U+\lfloor(\card V-\card U)/2\rfloor$. By the
Tutte--Berge formula, the minimum of the right side over all $U$, with the components of $G-U$ as blocks, is the matching
number $\nu(G)$. Conversely, over every field and for every graph, a matching of $m$ edges gives such a matrix with
$\dim\ker M\ge m$. Hence, for triangle-free $G$ of maximum degree at most~$3$, $\minrk_F(G)=\card V-\nu(G)$ over every
field.
\end{theorem}

\begin{outline}
Deleting $U$ costs at most $\card U$ dimensions. A code on a union of blocks with no edges between them splits as a sum
of the codes on the blocks, because the rows of one block vanish on the others; each block's code obeys \cref{thm:S}. The
matching code puts one free symbol on each edge of the matching, constant along the edge, and zero elsewhere; it is a
linear storage code of dimension~$m$, and \cref{lem:linear-pattern} turns it into a matrix with $\dim\ker M\ge m$. The
Tutte--Berge formula is classical and not formalized.
\end{outline}

The statement appears to be new: we found no statement that the minrank of triangle-free subcubic graphs equals the
clique-cover number $\card V-\nu(G)$ over every field, in the minrank, index-coding, storage-code or guessing-game
literature. Special cases are known. Tims~\cite{Tims2013} proves that the minimum rank over
every field equals the clique-cover number for all graphs of order at most~$10$ (by computer), the Petersen graph
included, and for the Möbius ladders, cubic triangle-free graphs for which Blasiak, Kleinberg and
Lubetzky~\cite{BKL2013} had shown broadcast rate $\card V/2$ for all codes; Berliner and
Langberg~\cite{BerlinerLangberg2011} prove the same equality over every finite field for outerplanar graphs. None of the hypotheses can be dropped.
The Clebsch graph (degree~$5$, \cref{cor:H4}) has binary minrank $6<8=\card V-\nu$. The complete graph $K_4$ shows that
triangles must be excluded. On $C_5$, vector-linear codes with two symbols per vertex reach dimension~$5$, more than
$2\nu(C_5)=4$, so the theorem is genuinely scalar (although rate~$\tfrac12$ still holds for $C_5$ by \cref{prop:C5}).

\paragraph{Why it matters.} Theorem~S settles, for scalar linear codes, the first whole class of bounded
degree beyond degree~$2$, where earlier results covered particular families, and it is exact: in degree~$3$ no scalar
linear code over any field, characteristic~$2$ included, beats the pairing on any triangle-free graph, locally
tree-like or not. With Theorem~A it narrows where a coding gain could live: for scalar linear codes, at degree at
least~$4$ over a finite field through $p$-torsion; for vector-linear codes over finite fields and for nonlinear codes,
already at degree~$3$.
